\documentclass{article}
\usepackage{amsmath}
\usepackage{amssymb}
\usepackage{mathtools}
\usepackage{amsthm}
\usepackage{bbm}
\usepackage{hyperref}

\title{Jieba dictionary training derivation}
\author{Kaiwen Wu}
\date{October 12, 2026}

\newcommand{\jb}{\texttt{jieba}}

\newtheorem{theorem}{Theorem}

\begin{document}
\maketitle

\section{Preliminaries and problem setup}
\jb\ (\url{https://github.com/fxsjy/jieba}) is a well-known Chinese word segmentation library in Python.
It comes with a default dictionary $D$ of Chinese words $w$, each paired with a frequency score $f_w$.
At inference time, it builds a directed acyclic graph (DAG) of dictionary words from every sentence.
The edge $\mathsf{e}_{i \to j}$ in the DAG represents a dictionary word spanning from character $i$ (inclusive) to character $j$ (exclusive), with weight $s_\mathsf{e} \coloneqq \log f_\mathsf{e} - \log f_0$, where $f_0 = \sum_{w \in D} f_w$.
Notice that we abuse the edge notation to also mean the corresponding dictionary word.
Thus, for a sentence of length $m$, a path $\mathsf{p}$ the connects node 0 to node $m$ defines a segmentation, and the sum of the edge weights $s_\mathsf{p} \coloneqq \sum_{\mathsf{e} \in \mathsf{p}} s_\mathsf{e}$ is the unigram log likelihood of a sequence of length $|\mathsf{p}|$.\footnote{Strictly speaking, the edge weights sum is \emph{not} the log likelihood of the segmentation $\mathsf{p}$, as we'll explain in the next section.}
Since the topological ordering of the nodes is known, we're able to decode the best segmentation $\mathsf{p}_\ast$ with the largest $s_\mathsf{p}$ efficiently by dynamic programming (DP).
Some more details are:

\begin{itemize}
    \item Out-of-vocabulary (OOV) words are split into characters $\mathsf{e}_{i \to i+1}$ and assigned edge weight $\log{1}-\log f_0$ when building the DAG. Hence, if there are many OOV words, the edge weights can no longer be interpreted as normalized log-probs.
    \item A hidden Markov model is optionally applied to OOV fragments to provide finer segmentation.
\end{itemize}

Our task is to augment the default dictionary $D$ with a user dictionary $U$, and hope that the word segmentation quality is not impaired by the addition.
We cannot assume the persistence of quality because $U$ may contain frequency scores incompatible with $D$, e.g.\@ one that originates from \url{https://github.com/iDvel/rime-ice}, and thus creating competing candidate paths in the DAG that negatively affects the decision.
It's worth noting that word frequencies from $U$ are sometimes adjusted in an ad hoc manner (see \url{https://github.com/fxsjy/jieba/issues}), which is inefficient and error-prone.
Since $D \cup U$ is often sufficently large, we assume absence of OOV words thereafter to eliminate the cluttering of handling them.

\section{A generative unigram model of sentences}
Out of the DAG implementation, we prefer a different set of notations.
In this section, we use the following notations:\@
$S$ for a segmentation, which is a sequence of dictionary words;
$x$ for an unsegmented sentence;
$S_\oplus$ for the concatenation of the words in $S$, which itself is a sentence.
We denote by $D$ to refer the dictionary in general sense, since whether it's the default or the augmented dictionary usually does not matter.
We also denote the parameters of a language model by $\boldsymbol\phi$.
For example, $\boldsymbol\phi$ often includes the dictionary frequency scores $f_w$, whose normalized counterpart in terms of the dictionary is denoted by $\theta_w$.

\jb's segmentation log-prob $s_\mathsf{p}$, can be rewritten in our current parlance as $\log \prod_{w \in S} \theta_w$.
We note that $\prod_{w \in S} \theta_w$ is \emph{not} the likelihood $P(S; \theta)$ of observing $S$.
It's instead the condition probability $P(S \mid |S|\ \text{fixed}; \theta)$.

\begin{theorem}
    The sum of $P(S; \theta)$ over all possible $S$ does not sum to 1.
\end{theorem}
\begin{proof}
    \[
        \begin{aligned}
            \sum_S P(S; \theta)
            &= \sum_{m=0}^\infty \sum_{S: |S|=m} P(S; \theta)\\
            &= \sum_{m=0}^\infty \sum_{S: |S|=m} \prod_{w \in S} \theta_w\\
            &= \sum_{m=0}^\infty \prod_{i=1}^m \sum_{w_i \in D} \theta_{w_i}\\
            &= \sum_{m=0}^\infty \prod_{i=1}^m 1\,,
        \end{aligned}
    \]
    
    \noindent
    and the r.h.s.\@ diverges.
    Notice that as a byproduct we have shown that $\forall m$,

    \[
        P(S \mid |S|=m; \theta) = \prod_{w \in S} \theta_w\,.
    \]
\end{proof}

The reason is that there is no stopping rule in the generative story of \jb's likelihood.
For this purpose, there are two options:
(1) define $\hat{D} \coloneqq D \cup \{\epsilon\}$ and stop generating new words after $\epsilon$ has been sampled, then $\hat{\theta}_w \propto \theta_w$;
(2) for each token to be generated, first toss a biased coin whose head probability is $\tau$ to decide whether to stop, and then sample a token from $D$ if we have decided to continue.
It's straightforward to show that if $\tau = \hat{\theta}_\epsilon$, then $\hat{\theta}_w = \theta_w (1-\tau)$, or in other words, the two stories are equivalent.
Considering the fact that we ultimately need to return the unigram frequencies as output, parameterizing the stopping rule as an extra $\tau$ is more convenient, since it does not affect the normalization of $\theta_w$'s, neither does it affect how we will enforce prior on $\theta_w$'s.

\subsection{Learning}
Given a corpus of sentences $C$, we will train $\boldsymbol\phi = (\theta, \tau)$ by MLE/MAP of the log marginal likelihood.
We will achieve this by EM.

The EM auxiliary function is

\[
    \begin{aligned}
        Q(\boldsymbol\phi \mid \boldsymbol\phi^\mathrm{old})
        &= \sum_{x \in C} \sum_S P(S \mid x; \boldsymbol\phi^\mathrm{old}) \log P(S, x; \boldsymbol\phi)\\
        &= \sum_{x \in C} \sum_S P(S \mid x; \boldsymbol\phi^\mathrm{old}) (\log P(S; \boldsymbol\phi) + \log P(x \mid S))\,.
    \end{aligned}
\]

\noindent
Since $P(x \mid S) = \mathbbm{1}[x = S_\oplus]$ is irrelevant of $\boldsymbol\phi$, and the posterior $P(S \mid x)$ is only supported on region where $x = S_\oplus$, we can safely drop this term.
Then,

\begin{equation}\label{eq:em-aux}
    Q(\boldsymbol\phi \mid \boldsymbol\phi^\mathrm{old})
    = \sum_{x \in C} \sum_S P(S \mid x; \boldsymbol\phi^\mathrm{old}) \log P(S; \boldsymbol\phi)\,.
\end{equation}

\subsubsection{Learning \texorpdfstring{$\theta$}{theta}}
First, write:

\[
    \begin{aligned}
        \log P(S; \boldsymbol\phi)
        &= \log\left\{\tau (1-\tau)^{|S|} \prod_{w \in S} \theta_w\right\}\\
        &= \log\left\{\tau \prod_{w \in S} (1-\tau) \theta_w\right\}\\
        &= \log\left\{\tau \prod_{w \in S} \hat{\theta}_w\right\}\\
        &= \log\tau + \sum_{w \in S} \log \hat{\theta}_w\\
        &= \log\tau + \sum_{w \in D} N_w(S) \log \hat{\theta}_w\,,
    \end{aligned}
\]

\noindent
where $N_w(S)$ denotes the number of occurrences of $w$ in $S$.
Plugging in to Equation~\eqref{eq:em-aux} gives:

\[
    \begin{aligned}
        Q(\boldsymbol\phi \mid \boldsymbol\phi^\mathrm{old})
        &= \sum_{x \in C} \sum_S P(S \mid x; \boldsymbol\phi^\mathrm{old})\left(\log\tau + \sum_{w \in D} N_w(S) \log \hat{\theta}_w\right)\\
        &= \sum_{x \in C}\sum_S P(S \mid x; \boldsymbol\phi^\mathrm{old}) \sum_{w \in D} N_w(S) \log \hat{\theta}_w + \text{constant}\\
        &= \sum_{w \in D} \left[\sum_{x \in C} \sum_S P(S \mid x; \boldsymbol\phi^\mathrm{old}) N_w(S)\right] \log ((1-\tau) \theta_w) + \text{constant}\\
        &= \sum_{w \in D} \left[\sum_{x \in C} \sum_S P(S \mid x; \boldsymbol\phi^\mathrm{old}) N_w(S)\right] \log \theta_w + \text{constant}\,.
    \end{aligned}
\]

\noindent
Hence,

\[
    \theta_w \propto \sum_{x \in C} \mathbb{E}[N_w(S) \mid x; \boldsymbol\phi^\mathrm{old}]\,,
\]

\noindent
up to some additive offset from Dirichlet prior or pseudocounts.

\subsubsection{Learning \texorpdfstring{$\tau$}{tau}}
Again, from Equation~\eqref{eq:em-aux},

\[
    \begin{aligned}
        Q(\boldsymbol\phi \mid \boldsymbol\phi^\mathrm{old})
        &= \sum_{x \in C} \sum_S P(S \mid x; \boldsymbol\phi^\mathrm{old}) \log \left(\tau (1-\tau)^{|S|} \prod_{w \in S} \theta_w\right)\\
        &= \sum_{x \in C} \sum_S P(S \mid x; \boldsymbol\phi^\mathrm{old}) (\log\tau + |S| \log(1-\tau)) + \text{constant}\\
        &= |C| \log\tau + \left[\sum_{x \in C} \sum_S P(S \mid x; \boldsymbol\phi^\mathrm{old}) |S|\right] \log(1-\tau) + \text{constant}\,.
    \end{aligned}
\]

\noindent
Taking derivative w.r.t.\@ $\tau$ and making it zero yields:

\[
    \tau = \frac{1}{1 + \tilde{m}}\,,
\]

\noindent
where

\[
    \tilde{m} = \frac{1}{|C|} \sum_{x \in C} \mathbb{E}[|S| \mid x; \boldsymbol\phi^\mathrm{old}]\,.
\]

\subsubsection{Finding the expectation in \texorpdfstring{$\theta$}{theta} update}
To evaluate $\mathbb{E}[N_w(S) \mid x]$, first notice that for $S$ such that $S_\oplus = x$,

\[
    N_w(S)
    = \sum_{0 \le i < j \le |x|} \mathbbm{1}[x_{i:j} = w \land (i,j) \in S]
    = \sum_{(i,j): x_{i:j} = w} \mathbbm{1}[(i,j) \in S]
\]

\noindent
Hence,

\[
    \begin{aligned}
        \mathbb{E}[N_w(S) \mid x]
        &= \sum_S P(S \mid x) \sum_{(i,j): x_{i:j} = w} \mathbbm{1}[(i,j) \in S]\\
        &= \sum_{(i,j): x_{i:j} = w} \sum_{S: (i,j) \in S} P(S) \Big/ P(x)\\
        &= \sum_{(i,j): x_{i:j} = w} \sum_{S_p: (\cdot, i) \in S_p} \sum_{S_s: (j, \cdot) \in S_s} P(S_p \oplus (x_{i:j}) \oplus S_s) \Big/ P(x)\\
        &= \sum_{(i,j): x_{i:j} = w} \sum_{S_p: (\cdot, i) \in S_p} \sum_{S_s: (j, \cdot) \in S_s} \tau \left(\prod_{w_1 \in S_p} \hat{\theta}_{w_1} \cdot \hat{\theta}_w \cdot \prod_{w_2 \in S_s} \hat{\theta}_{w_2}\right) \Big/ P(x)\\
        &= \tau \sum_{(i,j): x_{i:j} = w} \hat{\theta}_w \sum_{S_p: (\cdot, i) \in S_p} \prod_{w_1 \in S_p} \hat{\theta}_{w_1} \sum_{S_s: (j, \cdot) \in S_s} \prod_{w_2 \in S_s} \hat{\theta}_{w_2} \Big/ P(x)\,,
    \end{aligned}
\]

\noindent
where

\[
    P(x) = \tau \sum_{S: (0, \cdot) \in S \land (\cdot, |x|) \in S} \prod_{w \in S} \hat{\theta}_w\,.
\]

\noindent
Here, $S_p$ denotes prefix segmentations that also satisifies $(0, \cdot) \in S_p$, and $S_s$ denotes suffix segmentation that also satisfies $(\cdot, |x|) \in S_s$.
Thus, this whole problem reduces to finding:

\[
    \tau^{-1} P(x_{s:t}) = \sum_{S: (s, \cdot) \in S \land (\cdot, t) \in S} \prod_{w \in S} \hat{\theta}_w\,.
\]

\noindent
and without loss of generality, let $s \coloneqq 0$ and $t \coloneqq n$.
We may find $\tau^{-1} P(x_{0:n})$ by DP.
Define $A_i \coloneqq \tau^{-1} P(x_{0:i})$.
We need to find $A_n$.
Then,

\[
    A_i \hat{\theta}_{x_{i:j}} = \sum_{S:(0,\cdot) \in S \land (\cdot, i) \in S} \prod_{w \in S \oplus (x_{i:j})} \hat{\theta}_w
\]

\noindent
And thus,

\[
    A_j = \sum_{i < j} A_i \hat{\theta}_{x_{i:j}}
\]

\subsubsection{Finding the expectation in \texorpdfstring{$\tau$}{tau} update}
To evaluate

\[
    \begin{aligned}
        \mathbb{E}[|S| \mid x]
        &= \sum_{S: (0,\cdot) \in S \land (\cdot,|x|) \in S} P(S \mid x) |S|\\
        &= \sum_{S: (0,\cdot) \in S \land (\cdot,|x|) \in S} P(S) |S| \Big/ P(x)\,.
    \end{aligned}
\]

\noindent
As before, we define $A_i \coloneqq \tau^{-1} P(x_{0:i})$.
We also define

\[
    B_i \coloneqq \tau^{-1} \sum_{S:(0,\cdot) \in S \land (\cdot,i) \in S} P(S) |S|\,.
\]

\noindent
We need to find $B_{|x|}$.
Then,

\[
    \begin{aligned}
        A_i \hat{\theta}_{x_{i:j}} + B_i \hat{\theta}_{x_{i:j}}
        &= \tau^{-1} \sum_{S:(0,\cdot) \in S \land (\cdot,i) \in S} \prod_{w \in S \oplus (x_{i:j})} \hat{\theta}_w \\
        &\quad\ + \tau^{-1} \sum_{S:(0,\cdot) \in S \land (\cdot,i) \in S} |S| \prod_{w \in S \oplus (x_{i:j})} \hat{\theta}_w\\
        &= \tau^{-1} \sum_{S:(0,\cdot) \in S \land (\cdot,i) \in S} (|S|+1) \prod_{w \in S \oplus (x_{i:j})} \hat{\theta}_w\\
        &= \tau^{-1} \sum_{S:(0,\cdot) \in S \land (\cdot,i) \in S} |S \oplus (x_{i:j})| \prod_{w \in S \oplus (x_{i:j})} \hat{\theta}_w
    \end{aligned}
\]

\noindent
Hence, $B_j = \sum_{i<j}(A_i + B_i)\hat{\theta}_{x_{i:j}}$.

\end{document}
